% =========================================================
% DU/TA/AM Framework
% =========================================================
\paragraph{The DU/TA/AM Framework.}

Every line of an algebraic proof is justified by exactly one of three
sources. This framework makes that obligation explicit and provides
the tagging system used throughout these notes.

\begin{tcolorbox}[colback=propbox, colframe=propborder, arc=2pt,
  left=6pt, right=6pt, top=4pt, bottom=4pt,
  title={\small\textbf{The Three Line Tags}}, fonttitle=\small\bfseries]
\begin{itemize}[itemsep=4pt]
  \item \tagDU\ \textbf{(Definition / Unpack).}
    The step applies a definition or unpacks a named concept into its
    explicit conditions.
    \textit{Example}: expanding ``$e$ is an identity'' into
    ``$ea = ae = a$ for all $a \in G$''.

  \item \tagTA\ \textbf{(Theorem / Axiom).}
    The step cites a previously proved proposition or a structural
    axiom (group axiom, ring axiom, field axiom, Peano axiom).
    \textit{Example}: citing G1 (associativity) to regroup $(ab)c = a(bc)$.

  \item \tagTA\ \textbf{(Algebraic Manipulation, abbreviated AM).}
    The step performs a computation whose justification is itself a
    combination of DU and TA moves, compressed for readability.
    \textit{Example}: simplifying $a \cdot e = a$ in one step after
    the identity axiom has already been cited.
\end{itemize}
\end{tcolorbox}

\begin{remark}[Why the framework matters: finding where proofs break]
The DU/TA/AM system serves a diagnostic purpose. Every unjustified step
in a proof is either a DU step that has not been taken (some definition
has not been unpacked) or a TA step that has not been cited (some axiom
or theorem is being used silently).

When an algebraic proof stalls, the right question is: which definition
have I not yet unpacked? Which axiom am I using but not citing? The
tags make these omissions visible.

In practice, stalling almost always traces to a definition that has
been left in its named form (``$G$ is a group'') when it needed to be
expanded into operative conditions (``associativity holds; an identity
exists; every element has an inverse''). Expand the definition and
the stuck proof typically becomes unstuck.
\end{remark}

\begin{remark}[The hierarchy: DU and TA are atomic; AM is shorthand]
DU and TA are the atomic justifications. There is nothing more basic.
AM is not a separate category of reasoning --- it is a compressed
sequence of DU and TA steps that have been verified and labeled as
a single unit for concision.

A fully formal proof contains only DU and TA steps, explicitly tagged.
The three-column proof format used in this project --- tag in the left
column, statement in the middle, commentary in the right --- enforces
explicit justification of every line and is the recommended format
while you are learning.

As fluency develops, AM steps can absorb more and more of the lower-level
manipulations. But the underlying DU/TA structure is always present,
even when not written. The ability to decompose any AM step into its
DU and TA components is the hallmark of genuine understanding.
\end{remark}

\begin{remark}[A proof is a chain of citations]
The deepest way to understand the DU/TA/AM system is this: a mathematical
proof is a sequence of statements, each of which is justified by citing
something that is already known (a definition, an axiom, or a previously
proved theorem). The proof is valid if and only if every step has a valid
citation.

When you read a proof and feel uncertain whether it is correct, the
right move is to label every step. If you can assign a DU or TA
justification to every line, the proof is valid. If any line has no
justification, there is a gap.

This is not just a learning technique. Working mathematicians who find
a suspicious step in a proof do exactly this: they ask ``what justifies
this?'' and try to answer. If no answer is available, the step is an error.
\end{remark}
